\origpage{25}{}
Let $F, F'$ be good lines with $\Delta(F') \subseteq \Delta(F)$. Let $F$ have ends $u_1, u_2$, and let $F'$ have ends $u'_1, u'_2$, so that $u_1, u'_1, u'_2, u_2$ are in order in $F^*$.

\hspace*{12pt}(2) \textit{$A(F) \subseteq A(F')$ and $B(F') \subseteq B(F)$.}

This is immediate since $\Delta(F') \subseteq \Delta(F)$.

Let $I_1 \subseteq F^*$ be the support with ends $u_1, u'_1$, and let $I_2 \subseteq F^*$ have ends $u'_2, u_2$. From (6.2), there are $q_1$ mutually vertex-disjoint paths of $\mu(I_1)$ between $V(\mu(u_1))$ and $V(\mu(u'_1))$, and similarly for $I_2$, where the presentation has depth $q_1$. From (6.1)(iii), we deduce that

\hspace*{12pt}(3) \textit{If $|X(F)| = |X(F')|$ there are $|X(F)|$ mutually vertex-disjoint paths of $B(F) \cap A(F')$ between $X(F)$ and $X(F')$, each using no edge of $\Gamma$.}

\hspace*{12pt}(4) \textit{If $F_1, \ldots, F_n$ is a sequence of good lines such that $\Delta(F_{i+1}) \subseteq \Delta(F_i)$ for $1 \le i < n$, there exists $i$ with $1 \le i < n$ such that every edge of $\Gamma \cap \Delta(F_i)$ is an edge of $\Gamma \cap \Delta(F_{i+1})$.}

\textit{Subproof.} From (1), (2), (3) and (2.6) (with $k$ replaced by $2q$), we deduce that there exists $i, i'$ with $1 \le i < i' \le n$, such that $|X_i| = |X_{i'}|$ and $L \cap B_i \cap A_{i'}$ uses no edge of $\Gamma$. But every edge of $B_i \cap A_{i'}$ which is in $\Gamma$ is also in $L$, since $L$ is exhaustive, and so $E(\Gamma \cap B_i \cap A_{i'}) = \emptyset$. Then every edge of $\Gamma \cap \Delta(F_i)$ is also an edge of $\Gamma \cap \Delta(F_{i'})$ and hence of $\Gamma \cap \Delta(F_{i+1})$. This proves (4).

Let us say that two distinct regions $r_1, r_2$ of $\Gamma$ in $\Sigma$ \textit{touch} if there is an edge $e$ of $\Gamma$ with $e \subseteq \bar{r}_1 \cap \bar{r}_2$. It follows that if $r_1, r_2$ touch then $\bar{r}_1 \cap \bar{r}_2 = U(J)$ for some $\Sigma$-jump $J$, and if they do not touch then $|\bar{r}_1 \cap \bar{r}_2| \le 1$. In particular, the touching relation defines a graph $T$ with vertex set the set of regions of $\Gamma$ in $\Sigma$, and it is a tree. From (4), every vertex of $T$ has degree at most $n$, and every path of $T$ \fix{with one end $r_0$, the region of $\Gamma$ that meets $bd(\Sigma) \setminus F^*$,} has at most $n$ vertices.\fixnote{fix:6.4-tree}{As printed, the claim is about every path of $T$, and (4) does not give it. Claim (4) is about nested discs, and the discs $\Delta(F)$ lie on the side of $F$ away from $r_0$. A path through $r_0$ joins two paths that start at $r_0$, so (4) bounds it only by $2n - 1$ vertices. For example, if $\Gamma$ consists of $k$ nested jumps on each of two disjoint intervals of $F^*$, then $T$ is a path with $2k + 1$ vertices and middle vertex $r_0$, and (4) bounds $k$, not $2k$. Paths that start at $r_0$ are enough for the count, as the inserted text shows.} \fix{For suppose that $r_0, r_1, \ldots, r_n$ is a path of $T$. For $1 \le i \le n$ let $P_i$ be the $\Sigma$-jump with $\bar{r}_{i-1} \cap \bar{r}_i = U(P_i)$, and let $F_i$ be a good line with the same ends as $P_i$ and with interior in $r_{i-1}$. Since $\Delta(F_i)$ is on the side of $F_i$ away from $r_0$, it follows that $\Delta(F_{i+1}) \subseteq \Delta(F_i)$, and the edges of $P_i$ are in $\Delta(F_i)$ but not in $\Delta(F_{i+1})$ $(1 \le i < n)$, contrary to (4). Consequently, every vertex of $T$ is at distance $< n$ from $r_0$. Since every vertex of $T$ has degree at most $n$, for $0 \le d < n$ there are at most $n^d$ vertices at distance $d$ from $r_0$, and so $|V(T)| \le 1 + n + \dots + n^{n-1}$.} Hence $|E(T)| \le n^n = \lambda$. But $E(T)$ is in 1-1 correspondence with the $\Sigma$-jumps, and the result follows. $\square$

\textbf{6.5} \textit{For all integers $p \ge 0$ and $q \ge 1$ there exists $\lambda \ge 0$ such that, if $G$ has a presentation in a disc $\Sigma$ of depth $\le q$ and defect 2, and there is an exhaustive vital $p$-linkage in $G$, then there are at most $\lambda$ $\Sigma$-jumps.}
